\section{Four-point potentials}\label{sec:green}

The conductivity theorem identifies local energies of uniformly convergent
functions.  Point-source voltages have no bounded total energy in the
limit, so we first obtain bounds away from their poles.  We then use local
energy replacements to determine their harmonic limit and its flux.

Throughout this section \(m=2n\), and the network is the full primal
network on the sphere.  Its vertex positions are \(z_n(v)\) from
Proposition~\ref{prop:geometry}.  Choose an edge position \(z_n(e)\) on
each canonical edge arc.  For a spatial Borel set \(U\), write
\[
 \mathcal E_n^0(w;U)
   =\sum_{e:\,z_n(e)\in U}(d_e w)^2.
\]
This is the notation \eqref{eq:local-network-energy} in the canonical
coordinate.  The choice of edge positions has no effect on inner limits.

Let \(\mathsf a_n,\mathsf b_n\) be distinct vertices.  The unit-current voltage is a
solution of
\begin{equation}\label{eq:point-current}
                   Q_n u_n=\delta_{\mathsf a_n}-\delta_{\mathsf b_n},
\end{equation}
where \(Q_n\) is the graph Laplacian of Lemma~\ref{lem:normalization}.
It exists and is unique up to a constant because the graph is connected.
For every vertex function \(w\), summation by parts gives
\begin{equation}\label{eq:point-current-weak}
 \mathcal E_n^0(u_n,w)=w(\mathsf a_n)-w(\mathsf b_n).
\end{equation}
Equivalently, \(u_n\) minimizes
\begin{equation}\label{eq:source-functional}
 \mathcal J_n(v)
     =\mathcal E_n^0(v)-2\bigl(v(\mathsf a_n)-v(\mathsf b_n)\bigr).
\end{equation}
Indeed, for any \(v\),
\[
 \mathcal J_n(v)-\mathcal J_n(u_n)
 =\mathcal E_n^0(v-u_n)
   +2\mathcal E_n^0(u_n,v-u_n)
   -2\bigl((v-u_n)(\mathsf a_n)-(v-u_n)(\mathsf b_n)\bigr)
 =\mathcal E_n^0(v-u_n).
\]
This identity fixes the normalization of the flux calculation below.

Identify the canonical Riemann sphere with the unit round sphere in
\(\R^3\), and write
\(\operatorname{ch}(z,w)=\abs{z-w}_{\R^3}\) for chordal distance.
For distinct \(x,y,z_0\), define on \(\Sph\setminus\{x,y\}\)
\begin{equation}\label{eq:round-voltage}
 U_{x,y;z_0}(z)
   =\frac1{2\pi}
       \log\frac{\operatorname{ch}(z,y)\operatorname{ch}(z_0,x)}
                     {\operatorname{ch}(z,x)\operatorname{ch}(z_0,y)}.
\end{equation}
It vanishes at \(z_0\).  The ratio is the modulus of a conformal cross
ratio, so the function does not depend on the choice of stereographic
coordinate.

\begin{theorem}[Four-point voltage convergence]\label{thm:four-point}
Fix a joint extraction on which
Proposition~\ref{prop:geometry}, Theorem~\ref{thm:network-bounds}, and
Theorem~\ref{thm:conductivity} hold.  Let
\(\mathsf a_n,\mathsf b_n,\mathsf c_n,\mathsf d_n\) be vertex
sequences whose canonical positions converge respectively to four
pairwise distinct points \(x,y,z,z_0\).  Normalize the solution of
\eqref{eq:point-current} by \(u_n(\mathsf d_n)=0\).  Then, for every compact
\(K\subset\Sph\setminus\{x,y\}\),
\begin{equation}\label{eq:local-voltage-convergence}
 \sup_{v:\,z_n(v)\in K}
    \abs{u_n(v)-U_{x,y;z_0}(z_n(v))}\longrightarrow0.
\end{equation}
An empty supremum in \eqref{eq:local-voltage-convergence} is assigned
the value zero.
In particular,
\[
 u_n(\mathsf c_n)-u_n(\mathsf d_n)\longrightarrow
 \frac1{2\pi}
       \log\frac{\operatorname{ch}(z,y)\operatorname{ch}(z_0,x)}
                     {\operatorname{ch}(z,x)\operatorname{ch}(z_0,y)}.
\]
The assertion holds for any such varying vertex sequences on the
extracted sample.
\end{theorem}

\subsection{Bounds away from the current sources}

\begin{lemma}[Local bounds for point-source voltages]\label{lem:voltage-bounds}
Under the assumptions of Theorem~\ref{thm:four-point}, let
\(O=\Sph\setminus\{x,y\}\).  For every compact \(K\subset O\) there
are finite sample-dependent constants \(M_K,M_K'\) such that, eventually,
\[
 \sup_{v:\,z_n(v)\in K}\abs{u_n(v)}\le M_K,
 \qquad
 \mathcal E_n^0(u_n;K)\le M_K'.
\]
Here the energy on \(K\) counts edges located in \(K\).
Every subsequence has a further subsequence converging locally uniformly
in vertex tests on \(O\) to a continuous function.
\end{lemma}

\begin{proof}
For a vertex \(v\), there is a path from \(v\) to \(\mathsf a_n\) on which the
voltage is at least \(u_n(v)\).  To prove this, let \(S\) be the component
of \(v\) in the induced superlevel graph
\(\{w:u_n(w)\ge u_n(v)\}\).  If \(\mathsf a_n\notin S\), connectivity supplies
a boundary edge, and every boundary difference from \(S\) to its
complement is strictly positive.  Consequently
\[
 \sum_{w\in S}(Q_nu_n)(w)
   =\sum_{\substack{e=\{w,w'\}\\w\in S,\ w'\notin S}}
                   (u_n(w)-u_n(w'))>0,
\]
whereas \eqref{eq:point-current} makes the left side nonpositive.
This is a contradiction.  The sublevel argument with signs reversed
gives a path from \(v\) to \(\mathsf b_n\) with voltage at most \(u_n(v)\).

For real \(\alpha<\beta\), let \(T\) clip a number to
\([\alpha,\beta]\), and put \(D=\beta-\alpha\).  Monotonicity and the
one-Lipschitz property of \(T\) give on every edge
\[
 (d_e Tu_n)^2\le d_e u_n\,d_e Tu_n.
\]
Using \eqref{eq:point-current-weak},
\begin{equation}\label{eq:clipped-energy}
 \mathcal E_n^0(Tu_n)
 \le \mathcal E_n^0(u_n,Tu_n)
 =Tu_n(\mathsf a_n)-Tu_n(\mathsf b_n)\le D.
\end{equation}

Fix concentric inner and outer coordinate disks and a circular collar
between them surrounding the inner disk, with the entire closed outer
disk contained in \(O\).  For large \(n\), both source vertices lie
outside the outer disk.
Choose any two vertices in the inner disk and let their voltages be
\(\alpha\le\beta\).  The superlevel path from the higher vertex to
\(\mathsf a_n\) and the sublevel path from the lower vertex to \(\mathsf b_n\) meet every
surrounding circuit in the collar.  The intersections are graph vertices,
because these paths and circuits belong to the same embedded network.
At the two intersections the clipped voltage takes the values \(\beta\)
and \(\alpha\).  Its total variation on each such circuit is at least \(D\).

Theorem~\ref{thm:network-bounds} supplies a circuit law with
\(\sum_e\tau(e)^2\le B\).  The expectation below is only over that
finite-network circuit law.  By \eqref{eq:clipped-energy},
\[
 D\le
 \EE\!\left[\sum_{e\in\gamma}\abs{d_e Tu_n}\right]
 =\sum_e\tau(e)\abs{d_e Tu_n}
 \le \sqrt{B\,\mathcal E_n^0(Tu_n)}
 \le\sqrt{BD}.
\]
Thus the voltage oscillation in the inner disk is at most \(B\),
independently of the total energy of \(u_n\).

The punctured sphere \(O\) is connected.  Cover a compact \(K\subset O\)
and finitely many paths joining it to \(z_0\) by finitely many such disks,
with consecutive disks overlapping.  Vertices occur in each overlap for
all large \(n\), by the vanishing mesh in
Proposition~\ref{prop:geometry}.  Chaining the oscillation bounds and using
\(u_n(\mathsf d_n)=0\), \(z_n(\mathsf d_n)\to z_0\), gives the asserted bound \(M_K\).
The chart cover in Proposition~\ref{prop:geometry}(iv) allows these disks
at every point of \(O\), including a point omitted by one rooting.  Under
the uniform conformal or anticonformal comparison, sufficiently small
fixed round collars contain and are contained in reference collars of
comparable relative shape.  Their fixed-test circuit bounds therefore
apply in the canonical coordinate.

Choose a slightly larger compact neighborhood \(K'\subset O\) of \(K\).
For large \(n\), both endpoints of every edge located in \(K\) lie in
\(K'\).  Clip \(u_n\) to an interval containing all its values on \(K'\).
This clipping leaves those edge differences unchanged.
\eqref{eq:clipped-energy} bounds their total energy by the finite length
of that interval, giving \(M_K'\).

For each fixed \(V\Subset W\Subset O\), both moving source vertices
are eventually outside \(\overline W\).  Equation
\eqref{eq:point-current} then makes \(u_n\) harmonic at every vertex
with full star in \(W\).  Apply Lemma~\ref{lem:harmonic-compactness}
on \(W\), using the local bounds just proved, to obtain compactness on
\(V\).  A countable exhaustion by such buffered pairs and a diagonal
extraction give the final subsequential compactness assertion on \(O\).
\end{proof}

\subsection{Local energy convergence and harmonicity}

Take a subsequence furnished by Lemma~\ref{lem:voltage-bounds}, and write
\(u\) for its locally uniform limit on \(O\).  Associate to its finite
voltages the edge energy measures
\[
             \xi_n=\sum_e(d_eu_n)^2\delta_{z_n(e)}.
\]
Their masses are bounded on every compact subset of \(O\).  A diagonal
vague extraction on this locally compact space gives a Radon measure
\(\xi\) on \(O\).

\begin{lemma}[No local energy defect]\label{lem:no-energy-defect}
On this extraction, \(u\in H^1_{\rm loc}(O)\), \(u\) is harmonic on \(O\),
and
\begin{equation}\label{eq:no-energy-defect}
                         \xi=\abs{\nabla u}^2\dd x
\end{equation}
in every conformal chart of \(O\).  Hence the edge energies converge on
every relatively compact continuity set of the measure on the right.
\end{lemma}

\begin{proof}
We first justify the lower measure bound with the necessary outer room.
Fix a compact \(K\) inside a relatively compact open \(D\subset O\).
Choose \(K\subset V\Subset W\Subset D\), with \(V,W\) open and
\(\xi(\partial W)=0\).  Such a \(W\) can be chosen among distance
neighborhoods with sufficiently small radius: the boundaries are disjoint
as the radius varies, and a finite measure gives positive mass to at most
countably many of them.  The local uniform convergence and the lower
bound of Lemma~\ref{lem:uniform-relaxation} give
\[
 \mathcal R(u,V)
 \le\liminf_n\mathcal E_n^0(u_n;W)=\xi(W)<\infty.
\]
The equality uses vague convergence on the relatively compact continuity
set \(W\).  These choices around each compact set and
Lemma~\ref{lem:energy-bounds} put \(u\) in \(H^1_{\rm loc}(O)\).
Theorem~\ref{thm:conductivity} now identifies
\(\mathcal R(u,V)=\int_V\abs{\nabla u}^2\).  Thus
\[
 \int_K\abs{\nabla u}^2\dd x\le\xi(D).
\]
Take the supremum over compact \(K\subset D\), and then over relatively
compact \(D\).  This proves
\begin{equation}\label{eq:energy-defect-lower}
                        \abs{\nabla u}^2\dd x\le\xi.
\end{equation}
The continuity set \(W\) is essential here: the open-set Portmanteau
inequality alone has the opposite direction for the desired comparison.

We next use local minimization to obtain the reverse bound and harmonicity.
Work in one Euclidean chart, and let \(B_r=B(z,r)\Subset O\) be a disk
with \(\xi(\partial B_r)=0\).  Let
\(\varphi\in C_c^\infty(B_r)\).  Choose fixed radii
\[
             r<r_1<r_2<r_3<r_4,\qquad B_{r_4}\Subset O.
\]
The function \(u+\varphi\) is continuous and in \(H^1\) on this
buffered disk.  The recovery clause of
Theorem~\ref{thm:conductivity} gives \(w_n\to u+\varphi\) uniformly on a
neighborhood of \(\overline B_{r_3}\) and
\begin{equation}\label{eq:local-recovery-bound}
 \limsup_n\mathcal E_n^0(w_n;B_{r_3})
 \le\int_{B_{r_4}}\abs{\nabla(u+\varphi)}^2\dd x.
\end{equation}
Choose a cutoff \(\theta_n\) equal to one on \(B_{r_1}\) and zero outside
\(B_{r_2}\), with bounded energy for these fixed radii.  Put
\[
                 v_n=\theta_n w_n+(1-\theta_n)u_n.
\]
Extend the recovery arbitrarily where it is multiplied by zero.
All changed edges lie in \(B_{r_3}\) for large \(n\), by the open room
and star mesh.  The transition collar is outside \(B_r\), where
\(\varphi=0\), so \(w_n-u_n\to0\) uniformly there.  Applying the
edge identity \eqref{eq:edge-gluing} to these two functions gives
\begin{equation}\label{eq:local-replacement-bound}
 \mathcal E_n^0(v_n;B_{r_3})
 \le \mathcal E_n^0(w_n;B_{r_3})
       +\mathcal E_n^0(u_n;B_{r_3}\setminus\overline B_r)+o(1).
\end{equation}
Indeed the principal squared differences are bounded by their convex
combination.  The coefficient of an old difference can be nonzero only
outside \(\overline B_r\), since \(\theta_n=1\) on \(B_{r_1}\) and
the mesh is small.  The remaining term has energy bounded by
\(\sup_{\rm collar}\abs{w_n-u_n}^2\mathcal E_n^0(\theta_n)=o(1)\);
its cross term is \(o(1)\) by the bounded old and recovered energies.

Neither source value changes.  The minimizing identity
\eqref{eq:source-functional}, restricted to the changed edges, gives
\(\mathcal E_n^0(u_n;B_{r_3})
 \le\mathcal E_n^0(v_n;B_{r_3})\).
Subtract the old energy on
\(B_{r_3}\setminus\overline B_r\) in
\eqref{eq:local-replacement-bound}.  The difference on the left is
\(\xi_n(\overline B_r)\ge\xi_n(B_r)\), so the result is
\[
                \xi_n(B_r)
                \le\mathcal E_n^0(w_n;B_{r_3})+o(1).
\]
Take \(n\) to infinity.  The left side tends to \(\xi(B_r)\) because
\(B_r\) is a continuity set, and
\eqref{eq:local-recovery-bound} bounds the right side.  Now decrease
\(r_4\) to \(r\), choosing the three intermediate radii afresh between
them at each step.  Each cutoff is used only for a fixed step before
the network limit.  The continuum energy has no mass on the circle
\(\partial B_r\), so
\begin{equation}\label{eq:local-minimality-limit}
       \xi(B_r)\le
         \int_{B_r}\abs{\nabla(u+\varphi)}^2\dd x.
\end{equation}

Set \(\varphi=0\).  Together with
\eqref{eq:energy-defect-lower}, this gives equality of the two measures
on every such disk.  These disks form a cover basis after excluding
at most countably many radii at each center.  The nonnegative measure
\(\xi-\abs{\nabla u}^2\dd x\) vanishes on that basis, hence vanishes
on \(O\).  This proves \eqref{eq:no-energy-defect}.

For \(\psi\in C_c^\infty(B_r)\), use
\(\varphi=t\psi\) in \eqref{eq:local-minimality-limit} and the measure
equality.  Expanding for both signs of \(t\) and letting \(t\to0\)
gives
\(\int_{B_r}\nabla u\cdot\nabla\psi\dd x=0\).
Every point has such a disk, so \(u\) is distributionally harmonic on
\(O\).  Interior harmonic regularity makes it classically harmonic.
Finally, vague convergence to the measure in
\eqref{eq:no-energy-defect} gives convergence on each of its relatively
compact continuity sets.
\end{proof}

The proof used recovery for the particular limiting function \(u\).
This is legitimate pathwise: Lemma~\ref{lem:uniform-relaxation} obtains
its statements for every continuous function from one countable cell
library.  No conditional kernel depending on the subsequently chosen
voltage limit is needed.

\subsection{Unit flux and the logarithmic poles}

\begin{proof}[Proof of Theorem~\ref{thm:four-point}]
Take any subsequence.  Lemma~\ref{lem:voltage-bounds} and a further
vague energy extraction give a locally uniform limit \(u\) on \(O\).
Lemma~\ref{lem:no-energy-defect} makes \(u\) harmonic there and identifies
its local energy measure.  The normalization \(u_n(\mathsf d_n)=0\) and local
uniform convergence imply \(u(z_0)=0\).

Choose a conformal disk about \(x\) whose closure avoids \(y\), with
\(x\) at its Euclidean center.  Let \(\chi\) be a smooth radial function
equal to one on a smaller disk and zero outside a larger disk, with
its transition compactly in the punctured disk.  Choose a closed-buffered
annular domain \(V\Subset O\) containing \(\supp\nabla\chi\) with
clearance from both boundary components.  Its two boundary circles
have zero \(\abs{\nabla u}^2\dd x\)-mass.

Fix \(t\in\R\).  Choose annuli \(V\Subset W\Subset U\Subset O\) by
slightly moving the two boundary radii, and recover \(u+t\chi\) on a
neighborhood of \(\overline W\).  The recovery \(w_n\) satisfies
\begin{equation}\label{eq:flux-recovery}
 \limsup_n\mathcal E_n^0(w_n;W)
       \le\int_U\abs{\nabla(u+t\chi)}^2\dd x.
\end{equation}
For a vertex of the full network put
\[
                  \widehat u_n(v)=u_n(v)+t\chi(z_n(v)).
\]
Take an annular cutoff \(\theta_n\) equal to one on a neighborhood of
\(\overline V\) and zero outside an intermediate annulus compactly
contained in \(W\).  Such a cutoff is obtained by combining two disk
cutoffs; its energy is bounded for these fixed radii.  Define
\[
                  v_n=\theta_n w_n+(1-\theta_n)\widehat u_n.
\]
On both transition collars \(w_n-\widehat u_n\to0\) uniformly.  Wherever
the coefficient \(1-\theta_n\) of an edge difference can be nonzero,
the edge lies outside \(\overline V\) with a margin, where \(\chi\)
is constant on each component.  Thus
\(d_e\widehat u_n=d_eu_n\) on every such edge for all large \(n\).
This uses the clearance from \(\supp\nabla\chi\) and vanishing star mesh.
Outside \(W\),
\(\chi\) is constant on the inner and outer components and
\(\theta_n=0\); hence the changed energy is confined to \(W\).
The source vertex is in the inner component and the sink is in the
outer one, so
\[
                v_n(\mathsf a_n)=u_n(\mathsf a_n)+t,\qquad
                v_n(\mathsf b_n)=u_n(\mathsf b_n).
\]

The edge gluing estimate, with the old part supported outside
\(\overline V\), gives
\[
 \mathcal E_n^0(v_n;W)
 \le \mathcal E_n^0(w_n;W)
       +\mathcal E_n^0(u_n;W\setminus\overline V)+o(1).
\]
The error has energy at most the squared uniform collar discrepancy times
the fixed cutoff energy, and its cross term tends to zero by local energy
bounds.  Minimality of the source functional
\eqref{eq:source-functional} now gives
\[
 \mathcal E_n^0(u_n;W)
       \le\mathcal E_n^0(v_n;W)-2t.
\]
Subtract the displayed old collar energy.  The difference on the left is
\(\xi_n(\overline V)\ge\xi_n(V)\), so the resulting inequality is
\[
                  \xi_n(V)
                  \le\mathcal E_n^0(w_n;W)-2t+o(1).
\]
By Lemma~\ref{lem:no-energy-defect}, the left side converges to
\(\int_V\abs{\nabla u}^2\), since \(V\) is a continuity set.
Use \eqref{eq:flux-recovery}, then shrink \(U\) to \(V\) through annuli
with the same clearance from \(\supp\nabla\chi\).  At each stage the
annuli and cutoffs are fixed before \(n\) tends to infinity.  Absolute
continuity of the continuum energy and the zero mass of the boundary
circles give
\begin{equation}\label{eq:flux-variation}
 \int_V\abs{\nabla u}^2\dd x
 \le\int_V\abs{\nabla(u+t\chi)}^2\dd x-2t.
\end{equation}
Expanding this inequality yields
\[
 0\le
 2t\left(\int_V\nabla u\cdot\nabla\chi\dd x-1\right)
       +t^2\int_V\abs{\nabla\chi}^2\dd x.
\]
It holds for every real \(t\), so the coefficient of \(t\) is zero:
\begin{equation}\label{eq:unit-source-flux}
                     \int\nabla u\cdot\nabla\chi\dd x=1.
\end{equation}
For a cutoff equal to one near \(y\) and zero near \(x\), the replacement
changes the sink value by \(t\) and the source value by zero.
The source term in \eqref{eq:source-functional} then changes by \(+2t\);
the expanded inequality gives flux \(-1\) at \(y\).

We verify that these fluxes determine the only possible singularities.
Choose an outer coordinate disk about \(x\) whose closure excludes \(y\),
and a fixed collar inside it surrounding a smaller disk about \(x\).
For every sufficiently large \(n\), take a surrounding network circuit
in that collar.  Eventually
\(\mathsf a_n\) is in the smaller disk and \(\mathsf b_n\) is outside
the outer disk, hence outside the circuit.
Lemma~\ref{lem:voltage-bounds} bounds the voltages on that collar
uniformly.  Inside the circuit, \(Q_nu_n=0\) except at the positive
source, where \(Q_nu_n=1\).  The minimum principle therefore
bounds every interior voltage from below by the minimum on the circuit.
Passing to the limit bounds \(u\) from below near \(x\).  For \(y\),
choose an outer disk whose closure excludes \(x\) and a collar inside
it surrounding a smaller disk about \(y\), and take a surrounding
network circuit there.  Eventually \(\mathsf b_n\)
is inside its circuit and \(\mathsf a_n\) is outside the outer disk.
The maximum principle then bounds \(u\) from above near \(y\).

We recall the planar isolated-singularity argument at the needed strength.
If \(v\) is a positive harmonic function on a punctured disk, its circular
mean \(M(r)\) satisfies \((rM'(r))'=0\).  Thus
\[
                       M(r)=a\log(1/r)+b,\qquad a\ge0,
\]
where positivity forces the sign of \(a\).  Harnack's inequality on
the fixed-shape annulus \(\{r/2<\abs z<2r\}\) bounds the values on
\(\abs z=r\) by a constant times \(M(r)\), uniformly in small \(r\).
This version of Harnack follows directly from the Poisson kernel on a
finite chain of disks in the annulus.  Hence \(v=O(1+\log(1/r))\).
For each nonzero angular Fourier mode, harmonicity makes its radial
coefficient a linear combination of \(r^j\) and \(r^{-j}\),
\(j\ge1\).  The logarithmic bound forces every \(r^{-j}\) coefficient
to vanish.  The remaining positive-power series is the harmonic Poisson
extension from any fixed smaller circle.  Consequently
\[
                     v(z)=a\log(1/\abs z)+h(z),
\]
where \(h\) extends harmonically through the center.  Applying this to
\(u\) plus a constant near \(x\), and to \(-u\) plus a constant near
\(y\), shows that \(u\) has only logarithmic or removable singularities.

For \(a\log(1/\abs z)\), a radial cutoff equal to one near zero and zero
outside has
\[
 \int\nabla\!\bigl(a\log(1/\abs z)\bigr)\cdot\nabla\chi\dd x
       =-2\pi a\int\chi'(r)\dd r=2\pi a.
\]
The harmonic remainder contributes zero.  Equation
\eqref{eq:unit-source-flux} therefore sets the coefficient at \(x\)
to \(1/(2\pi)\), and the sink flux sets the coefficient at \(y\) to
\(-1/(2\pi)\).

For comparison, the round zero-mean Green kernel has, up to its additive
constant,
\[
                  G_{\rm rd}(z,w)
                  =-\frac1{2\pi}\log\operatorname{ch}(z,w).
\]
Away from \(w\), a radial calculation on the unit sphere gives
\(-\Delta_z G_{\rm rd}(z,w)=-1/(4\pi)\); the preceding local flux is
one at \(w\).  Thus its distributional equation is
\(-\Delta_zG_{\rm rd}(z,w)=\delta_w-(4\pi)^{-1}\dd\mathrm{vol}_{\rm rd}\).
The difference \(G_{\rm rd}(\cdot,x)-G_{\rm rd}(\cdot,y)\) has exactly
the two fluxes and singular coefficients just found.  Subtract it from
\(u\).  The difference extends harmonically across \(x,y\), and a
harmonic function on the compact sphere is constant.  Its value at
\(z_0\) gives precisely \eqref{eq:round-voltage}.

Every subsequence has therefore a further locally uniform subsequence with
the same limit \(U_{x,y;z_0}\).  This proves
\eqref{eq:local-voltage-convergence} along the entire extraction.
Evaluating at \(\mathsf c_n\) gives the asserted four-point value.  The reasoning
applied to arbitrary varying pole and evaluation vertices with the stated
distinct limits, so it proves the claimed pathwise quantifier.
\end{proof}
